\section{A faithful encoding by iterated integrals}
\label{sec:iterated-integrals}

We first encode noncommutative polynomials as power series in one variable.
The encoding will preserve every nonzero polynomial; the remaining sections
will bound how many Taylor coefficients are needed for a small formula.

Fix an integer $n\geq 1$ and a field $F$ of characteristic zero.
For $i\in[n]=\{1,\ldots,n\}$, define an $F$-linear operator on $F[[z]]$ by
\begin{equation}
  \mathcal J_i(a)(z)=\int_0^z\frac{a(t)}{t+i}\,dt.
  \label{eq:integration-operator}
\end{equation}
Here integration means termwise formal integration with zero constant term:
expand $(t+i)^{-1}=\sum_{j\geq0}(-1)^j t^j/i^{j+1}$ first.
All denominators are nonzero in $F$.
For a nonzero series $a$, let $\operatorname{ord}_0 a$ denote the least
exponent with nonzero coefficient, and put $\operatorname{ord}_0 0=\infty$.
Then
\begin{equation}
  \operatorname{ord}_0\mathcal J_i(a)
  \geq \operatorname{ord}_0 a+1.
  \label{eq:integration-order}
\end{equation}

A word $w=i_1\cdots i_b$ on $[n]$ has length $|w|=b$;
the empty word is denoted by $\epsilon$.
Define the series $h_w\in\mathbb Q[[z]]$ recursively by
\begin{equation}
  h_\epsilon=1,\qquad h_{iw}=\mathcal J_i(h_w).
  \label{eq:word-series}
\end{equation}
We regard these series as elements of $F[[z]]$ through the canonical
embedding $\mathbb Q\hookrightarrow F$.
In particular, $\operatorname{ord}_0 h_w\geq |w|$.
Over $\mathbb C$, they are holomorphic germs at zero: the same recursion
uses ordinary integration on any disk of radius less than one.
These functions are classical iterated integrals, also called
hyperlogarithms; see Chen~\cite{Chen1977} for the general framework.
The following independence statement is a special case of
Deneufch\^atel et al.~\cite[Theorem~1]{DeneufchatelEtAl2011}.
Their criterion in fact gives independence over $F(z)$: no nonzero
$F$-linear combination of $1/(z+1),\ldots,1/(z+n)$ is the derivative
of a rational function, because such a derivative has zero residue
at every pole. We give a direct proof of the constant-coefficient
independence needed here.
We use a triangular monodromy argument, as in the two-letter case treated
by Minh, Petitot and van der Hoeven
\cite[Section~5, Theorem~4]{MinhPetitotVanderHoeven1998}.

\begin{lemma}[Linear independence of the word series]
\label{lem:word-independence}
For every integer $n\geq1$, every integer $D\geq0$, and every
characteristic-zero field $F$, the series
\[
  \{h_w: w\text{ is a word on }[n],\ |w|\leq D\}
\]
are linearly independent over $F$.
\end{lemma}

\begin{proof}
The case $D=0$ consists of the single series $1$.
Suppose $D\geq1$, and first work over $\mathbb C$.
Let $E_D$ be the complex algebra with basis $X_w$, $|w|\leq D$,
whose multiplication is concatenation of words, with products of length
greater than $D$ set equal to zero.
Its identity is $X_\epsilon=1$.
The span $\mathfrak m$ of the nonempty words is an ideal satisfying
$\mathfrak m^{D+1}=0$.
For one-letter words write $X_i=X_{(i)}$, and set
\[
  U(z)=\sum_{|w|\leq D}h_w(z)X_w.
\]
The recursion \eqref{eq:word-series} gives the left differential equation
\begin{equation}
  U'(z)=\left(\sum_{i=1}^n\frac{X_i}{z+i}\right)U(z),
  \qquad U(0)=1.
  \label{eq:word-differential-equation}
\end{equation}
This is a linear equation on the finite-dimensional vector space $E_D$,
with coefficients holomorphic on
$\Omega=\mathbb C\setminus\{-1,\ldots,-n\}$.
For this particular equation, local existence and uniqueness follow
directly by induction on word length: the derivative of a length-$b$
coefficient is determined by length-$(b-1)$ coefficients. On a disk
avoiding the poles, taking holomorphic primitives with prescribed initial
values therefore constructs the solution in $D$ steps.
A finite chain of overlapping disks continues it along any path in
$\Omega$, and uniqueness makes the solutions agree on overlaps.

For a loop $\gamma$ in $\Omega$ based at zero, let $Q_\gamma$ be the
value at zero after continuation around $\gamma$.
It belongs to $1+\mathfrak m$, since the constant coefficient of the
solution remains one.
Suppose that $\sum_{|w|\leq D}c_w h_w=0$ as a complex germ, and define
a complex-linear functional $\lambda:E_D\to\mathbb C$ by
$\lambda(X_w)=c_w$.
The germ $\lambda(U)$ is zero, and its continuation around every loop
is zero; hence $\lambda(Q_\gamma)=0$ for every $\gamma$.
It therefore suffices to prove that the loop values span $E_D$.

The order of multiplication matters here.
Along a fixed path, a solution with initial value $b\in E_D$ equals
the solution with initial value one multiplied by $b$ on the right.
Thus traversing $\gamma$ and then $\delta$ gives
$Q_\delta Q_\gamma$.
In particular, the set of loop values contains $1$, and every ordered
product of loop values is again a loop value.

Choose a based loop $\gamma_i$ with winding number one about $-i$
and zero about every other pole, and put $Q_i=Q_{\gamma_i}$.
One may take an access path from zero, a small positive circle about
$-i$, and the reversed access path. For the degree-one integrals the
two access paths cancel. Direct integration on the circle gives
$2\pi\sqrt{-1}$ for $dz/(z+i)$ and zero for $dz/(z+j)$ when $j\ne i$;
the latter also follows by integrating its convergent Taylor expansion
on the disk bounded by the circle.
Reducing \eqref{eq:word-differential-equation} modulo $\mathfrak m^2$
and integrating around this loop yields
\[
  Q_i=1+\tau X_i\pmod{\mathfrak m^2},
  \qquad \tau=2\pi\sqrt{-1}.
\]
For $w=i_1\cdots i_b$, define
\[
  C_w=(Q_{i_1}-1)\cdots(Q_{i_b}-1),
  \qquad C_\epsilon=1.
\]
Expanding the product shows that $C_w$ lies in the complex linear span
of the loop values.
It has no terms of word length less than $b$, and its length-$b$ part
is $\tau^b X_w$.
Consequently the family $\{C_w:|w|\leq D\}$ is a basis of $E_D$:
in a word basis ordered by length, its coefficient matrix is triangular
with nonzero diagonal entries $\tau^{|w|}$.
The loop values therefore span $E_D$, so $\lambda=0$.
This proves complex linear independence.

To pass to $F$, form the matrix whose columns are the Taylor
coefficients of the finitely many series $h_w$, $|w|\leq D$.
Every entry is rational.
The kernels over $\mathbb C$ of its successive finite row prefixes
form a descending sequence of subspaces of a fixed finite-dimensional
space.
Their intersection is zero by the complex independence just proved,
so some finite prefix has full column rank.
It therefore contains a square minor with nonzero rational determinant.
That determinant stays nonzero under $\mathbb Q\hookrightarrow F$,
which rules out every $F$-linear relation among the series.
This minor is used only to prove independence; the construction of
the hitting point does not compute it.
\end{proof}

For $w=i_1\cdots i_b$, write $x_w=x_{i_1}\cdots x_{i_b}$ and
$x_\epsilon=1$.
Given $f=\sum_w f_w x_w\in F\langle x_1,\ldots,x_n\rangle$, define
\begin{equation}
  h_f=\sum_w f_w h_w.
  \label{eq:polynomial-series}
\end{equation}
Lemma~\ref{lem:word-independence} says that the $F$-linear map
$f\mapsto h_f$ is injective.
It is not an algebra homomorphism into $F[[z]]$.
Instead, the ordered word $x_{i_1}\cdots x_{i_b}$ acts by the composition
$\mathcal J_{i_1}\cdots\mathcal J_{i_b}$, applied to $1$:
\begin{equation}
  h_f=f(\mathcal J_1,\ldots,\mathcal J_n)1.
  \label{eq:operator-evaluation}
\end{equation}
For example, when $n\ge2$, the commutator $x_1x_2-x_2x_1$ gives
\[
 h_{12}(z)-h_{21}(z)=-\frac{z^3}{24}+O(z^4).
\]
Thus the encoding retains word order, but cancellation can delay the first
nonzero coefficient beyond the word length. We next use formula size to
bound this delay uniformly.
